\section{Casi reali: ricostruzione degli episodi}
\label{sec:episodi}

% Persone 2 e 3 — circa 1 pagina complessiva.
\subsection{chalk/debug, 8 settembre 2025}
% Persona 2 — circa 0,5 pagina.
\subsubsection{Cronologia essenziale e componenti coinvolti}

L'8 settembre 2025 l'account npm del manutentore fu usato per pubblicare
versioni sospette di pacchetti, tra cui \texttt{chalk@5.6.1} e
\texttt{debug@4.4.2}. Alle 13:16 UTC Aikido Security segnalò attività anomala
di pubblicazione; poco dopo il manutentore confermò di aver perso il controllo
dell'account. Segnalazioni sui repository GitHub di chalk e debug avvertirono
gli sviluppatori di non usare le versioni coinvolte~\cite{aikido2025,chalk2025}.

\begin{figure}[H]
    \centering
    \includegraphics[width=0.78\textwidth]{immagine_compromesso_chalk.png}
    \caption{Segnalazione della versione compromessa di chalk su GitHub.}
    {\footnotesize Fonte: \cite{chalk2025}\par}
    \label{fig:issue-chalk-compromesso}
\end{figure}

\subsubsection{Risposta adottata}

Le versioni compromesse furono rimosse da npm. L'advisory di debug raccomandò
agli utilizzatori di aggiornare alle versioni corrette e, se le versioni
malevole erano state installate, di rimuovere le dipendenze scaricate e
ricostruire da zero i bundle browser potenzialmente contaminati
~\cite{debug2025}. La cronologia qui riassunta descrive la risposta osservata;
le misure preventive generali sono trattate nel capitolo~\ref{sec:mitigazioni}.

\subsection{AUR, segnalazione del 12 giugno 2026}
\label{subsec:aur-episodio}
% Persona 3 — circa 0,5 pagina.
\subsubsection{Esempio e cronologia essenziale}
L'11 giugno 2026, sulla lista ufficiale \texttt{aur-general} fu segnalato un
aggiornamento sospetto del pacchetto AUR \texttt{gnome-randr-rust}: un nuovo
manutentore aveva aggiunto uno script che installava la dipendenza npm
\texttt{atomic-lockfile}. La segnalazione indicava che il pacchetto era stato
appena adottato e che l'aggiunta non era pertinente al software
~\cite{aur-mail-gnome-randr,aur-video-2026}. Il giorno successivo Arch pubblicò un avviso
relativo a un volume elevato di adozioni e aggiornamenti malevoli nell'AUR,
chiedendo agli utenti di esaminare le modifiche a \texttt{PKGBUILD} e script
d'installazione~\cite{arch-aur-incident-2026}. L'esempio mostra come una
modifica apparentemente ordinaria alla ricetta potesse introdurre un passaggio
verso npm, senza che il programma originale avesse bisogno di quella
dipendenza; i dettagli dell'hook e del suo funzionamento sono descritti nella
sezione~\ref{subsec:aur-tecnica}.

\subsubsection{Evidenze disponibili e risposta}
Nel thread, un manutentore AUR confermò che il commit era malevolo, ripristinò
la revisione precedente del pacchetto e sospese l'account che aveva pubblicato
la modifica~\cite{aur-mail-gnome-randr}. Questo documenta la risposta sul
repository, ma non prova che il pacchetto sia stato installato o che il codice
abbia raggiunto un sistema specifico. Nel successivo avviso, Arch spiegò che
stava rintracciando i commit malevoli e cercando di impedire nuove modifiche;
segnalò inoltre che, durante l'intervento, potevano verificarsi problemi nella
creazione di account, negli aggiornamenti e nelle adozioni di pacchetti
~\cite{arch-aur-incident-2026}. L'annuncio non offre una ricostruzione completa
del singolo caso né di tutti i pacchetti coinvolti; perciò l'esempio non va
esteso automaticamente all'intera ondata.
